Käesolevas peatükis esitatakse töö hetkeseis. Kuna eksperimendid ei ole veel lõppenud, on tegemist teadlikult vahekokkuvõttega, mitte lõplike järeldustega.

\section{Treeningutoru rekonstrueerimine}
Varasem \enquote{Kratt} treeningukatse ei andnud usaldusväärset tulemust. Hilisem sessioonilogide rekonstrueerimine näitas, et probleem ei olnud ainult mudeli kvaliteedis. Protsessis esines vähemalt kolm eri klassi vigu:
\begin{itemize}
    \item treeningukoodis ilmnes TensorFlow ja NumPy ühilduvusprobleem;
    \item evaluatsioonis puudus korrektselt moodustatud \texttt{testing\_ambient} kogum;
    \item andmestiku domeen ei kattunud reaalse seadme kasutusolukorraga.
\end{itemize}

Selle töö tulemusena parandati treeningu- ja raportiskripte ning muudeti evaluatsioonitugi selliseks, et ambient-andmete puudumine või vigane moodustamine oleks selgelt nähtav.

\section{Avaliku sanity-check eksperimendi tulemus}
Pärast kohaliku toru parandamist käivitati from-scratch sanity-check katse \texttt{Speech Commands} korpuse \texttt{marvin} sihtsõna peal \cite{speechcommands2018}. Selle katse eesmärk ei olnud lahendada eestikeelset ülesannet, vaid kontrollida, kas kohalik pipeline töötab otsast lõpuni avaliku kontrollitava sisendi peal.

Sanity-check eksperimendi tulemusena saadi edukalt:
\begin{itemize}
    \item tunnusefailid \texttt{mmap} vormingus;
    \item treenitud \texttt{microWakeWord} mudel;
    \item kvantiseeritud TFLite artefakt;
    \item automaatselt genereeritud analüüsikaust ROC-kõvera, threshold-tabelite ja score-jaotustega.
\end{itemize}

Selle tulemuse peamine väärtus seisneb selles, et treeningutoru enda töökindlus sai avaliku andmestiku peal kinnitust. Edasised eestikeelse mudeli probleemid ei ole seega enam automaatselt seletatavad infrastruktuuri veaga.

\section{Peamine metoodiline leid}
Senise töö üks olulisemaid järeldusi on seotud hindamisprotsessiga. Varasemas \enquote{Kratt} run'is saadud AUC väärtus ei osutunud mitte heale mudelile, vaid katkisele evaluatsioonile, sest \texttt{testing\_ambient} oli tühi. Pärast ambient-toe lisamist muutus ROC-analüüs sisuliseks ning oli võimalik eristada kahte olukorda:
\begin{itemize}
    \item degenerate tulemus, mis tuleb katkise sisendi tõttu;
    \item jämedateraline, kuid siiski korrektne tulemus, mis tuleb liiga lühikesest ambient-signaalist.
\end{itemize}

See tulemus on oluline, sest muudab töö rõhuasetust. Küsimus ei ole enam ainult selles, kas mudelit saab treenida, vaid ka selles, kas mudelit hinnatakse metodoloogiliselt korrektselt.

\section{Praegune seis}
Kirjutamise hetkeks on olemas:
\begin{itemize}
    \item toimiv kohalik treeningu-, ekspordi- ja analüüsitoru;
    \item avalikul andmestikul valideeritud \texttt{marvin} baseline;
    \item allalaadimisel ja ettevalmistamisel suuremad public ambient- ja speech-negative korpused;
    \item selgem arusaam sellest, kuidas tuleb üles ehitada eestikeelne \enquote{Kratt} andmestik.
\end{itemize}

Lõplikku tulemuste peatükki lisatakse hiljem koondtabelid valitud threshold'ide, \texttt{FRR}, \texttt{FAPH} ja seadmetestide kohta. Praeguses etapis on põhjendatud väita, et toru tehniline valideerimine avaliku andmestiku peal on õnnestunud.
